
\documentclass[12pt]{article}



\usepackage{amssymb}
\usepackage{amsmath}
\usepackage{ifthen}
\usepackage[table]{xcolor}
\usepackage{minitoc}
\usepackage{array}
\usepackage{graphicx}

\definecolor{yellow}{cmyk}{0,0,1,0}
\renewcommand{\arraystretch}{1.4}
\newcommand{\R}{\mathbb{R}}

\usepackage{colortbl}
\usepackage[top=0.75in, bottom=0.75in, left=1.25in, right=1.25in]{geometry}

%% Page size
%\setlength{\oddsidemargin}{13pt}
%%\setlength{\evensidemargin}{-0.5in}
%\setlength{\textheight}{674pt}
%\setlength{\textwidth}{426pt}
%\setlength{\topmargin}{-10pt}
%\setlength{\voffset}{-30pt}

\renewcommand{\arraycolsep}{3pt}


\input{color_flatex}


% Page size
%\setlength{\oddsidemargin}{-0.25in}
%\setlength{\evensidemargin}{-0.25in}
%\setlength{\textheight}{9.25in}
%\setlength{\textwidth}{7in}
%\setlength{\topmargin}{-0.5in}


% ----------------------------------------------------------------
\title{Triangular Matrix Matrix Multiplication \\ CS 378 Programming for Correctness and Performance}

\author{Devangi N. Parikh}

\begin{document}
	\maketitle

\section{Operation}
Consider the operation 

\[y := Ax + y,\]

where $A$ is a $m \times m$ matrix stored in the lower triangular part, and $x$
and $y$ are vectors of length $m$. This operation is a special case of
matrix-vector multiplication where $A$ is a symmetric matrix stored in the lower
triangular part. We will refer to this operation as {\sc Symv\_ln}m where the
{\sc ln} stands for 
\underline{l}ower
\underline{n}o-transpose.



\section{Precondition and Postcondition}

In the precondition 
\[
y = \widehat y
\]
$ \widehat y $ denotes the original contents of $ y $.
This allows us to express the state upon completion, the postcondition, as
\[
y = Ax + \widehat y
\]
It is implicitly assumed that $ A $  is symmetric and stored in the lower
triangular part. 

\section{Partitioned Matrix Expressions and Loop Invariants}

To derive the PME, partition 

\[ A \rightarrow \FlaTwoByTwo{A_{TL}}{A_{TR}}{A_{BR}}{A_{BL}},~
 y \rightarrow \FlaTwoByOne{y_{T}}{y_B}, \text{~and~} x \rightarrow \FlaTwoByOne{x_{T}}{x_B}  \]

Substituting these into the postcondition
yields
\[
\left(\begin{array}{c}
	y_T \\ \whline
	y_B
\end{array}\right)
=
\left(\begin{array}{c I c}
	A_{TL} & A_{BL}^T \\ \whline
	A_{BL} & A_{BR}
\end{array}\right) \left(
\begin{array}{c}
x_T \\ \whline x_B
\end{array}
\right) +
\left(\begin{array}{c}
	\widehat y_T \\ \whline
	\widehat y_B
\end{array}\right)
\]
or, equivalently,
\[
\left(\begin{array}{c}
	y_T \\ \whline
	y_B
\end{array}\right)
=
\left(\begin{array}{c}
	A_{TL}x_T + A_{BL}^Tx_B + \widehat y_T \\ \whline
	A_{BL}x_T + A_{BR}^Tx_B + \widehat y_B
\end{array}\right).
\]

From this, we can choose four loop invariants:
\begin{description}
	\item
	{\bf Invariant 1:}
\begin{equation}
	\left(\begin{array}{c}
		y_T \\ \whline
		y_B 
	\end{array}\right)  = 
	\left(\begin{array}{c}
		A_{TL} x_T + \widehat y_T \\ \whline 
		\widehat y_B 
	\end{array}\right)
\label{eq:inv1}
\end{equation}
	The top part has been partially computed and the bottom part has been left
alone. 
	\item
	{\bf Invariant 2:}
\begin{equation}
\left(\begin{array}{c}
		y_T \\ \whline
		y_B 
	\end{array}\right)  = 
	\left(\begin{array}{r}
		A_{TL} x_T + A_{BL}^T x_B + \widehat y_T \\ \whline 
		\widehat y_B 
	\end{array}\right)
	\label{eq:inv2}
	\end{equation} \\
	The top part has been fully computed and the bottom part has been left
alone.
\item 
{\bf Invariant 3:}
\begin{equation}
    \left(\begin{array}{c}
        y_T \\ \whline
        y_B 
    \end{array}\right)  = 
    \left(\begin{array}{c}
        A_{TL} x_T + \widehat y_T \\ \whline 
        A_{BL} x_T + \widehat y_B 
    \end{array}\right)
\label{eq:inv3}
\end{equation}
    The top part and bottom part has been partially computed.

\item
{\bf Invariant 4:}
\begin{equation}
    \left(\begin{array}{c}
        y_T \\ \whline
        y_B 
    \end{array}\right)  = 
    \left(\begin{array}{c}
        A_{TL} x_T +  A_{BL}^T x_B \widehat y_T \\ \whline 
        A_{BL} x_T + \widehat y_B 
    \end{array}\right)
\label{eq:inv3}
\end{equation}
 The top part is fully computed and bottom part has been partially computed.
 	\item
	{\bf Invariant 5:}
\begin{equation}
	\left(\begin{array}{c}
		y_T \\ \whline
		y_B 
	\end{array}\right)  = 
	\left(\begin{array}{c}
		 \widehat y_T \\ \whline 
		A_{BR}x_{B} + \widehat y_B 
	\end{array}\right)
\label{eq:inv5}
\end{equation}
	The top part has been partially computed and the bottom part has been left
alone. 
	\item
	{\bf Invariant 6:}
\begin{equation}
\left(\begin{array}{c}
		y_T \\ \whline
		y_B 
	\end{array}\right)  = 
	\left(\begin{array}{r}
		\widehat y_T \\ \whline 
		A_{BL}x_T + A_{BR}x_{B} + \widehat y_B 
	\end{array}\right)
	\label{eq:inv6}
	\end{equation} \\
	The top part has been fully computed and the bottom part has been left
alone.
\item 
{\bf Invariant 7:}
\begin{equation}
    \left(\begin{array}{c}
        y_T \\ \whline
        y_B 
    \end{array}\right)  = 
    \left(\begin{array}{c}
        A_{BL}^T x_T + \widehat y_T \\ \whline 
        A_{BR} x_B + \widehat y_B 
    \end{array}\right)
\label{eq:inv7}
\end{equation}
    The top part and bottom part has been partially computed.

\item
{\bf Invariant 8:}
\begin{equation}
    \left(\begin{array}{c}
        y_T \\ \whline
        y_B 
    \end{array}\right)  = 
    \left(\begin{array}{c}
        A_{BL}^T x_B+ \widehat y_T \\ \whline 
        A_{BL} x_T + A_{BR}x_B + \widehat y_B 
    \end{array}\right)
\label{eq:inv8}
\end{equation}
 The bottom part is fully computed and top part has been partially computed.
 
\end{description}

\section{Deriving the Algorithms}

\subsection{Loop Invariant \eqref{eq:inv1}}

The unblocked algorithm for loop invariant \eqref{eq:inv1} is showing in Figure~\ref{fig:unb_inv1}. While the blocked algorithm for loop invariant \eqref{eq:inv1} is showing in Figure~\ref{fig:blk_inv1}.


\resetsteps

\input{Derivations/trmm_llnn_unb_var1}

\begin{figure}
	\begin{center}
		\FlaWorksheet
	\end{center}
	\caption{Unblocked Algorithm for Loop Invariant \eqref{eq:inv1}}
	\label{fig:unb_inv1}
\end{figure}


\resetsteps

\input{Derivations/trmm_llnn_blk_var1}

\begin{figure}
	\begin{center}
		\FlaWorksheet
	\end{center}
	\caption{Blocked Algorithm for Loop Invariant \eqref{eq:inv1}}
	\label{fig:blk_inv1}
\end{figure}

\NoShow{
\section{Performance of the blocked and unblocked Algorithms}

Figure~\ref{fig:unbAlg1} shows the performance of the unblocked algorithm for loop invariant~\ref{eq:inv1} of TRMM.

\NoShow{
\begin{figure}
\includegraphics[width=\textwidth]{flamec/Plot_unb_GFLOPS.pdf}
\caption{Performance of the unblocked algorithm for loop invariant~\ref{eq:inv1} of TRMM}.
\label{fig:unbAlg1}
\end{figure}
}
}

\end{document}

